
The proliferation of pre-trained model repositories creates demand for automated architecture classification to enable model zoo management, compatibility verification, and architecture-aware model selection. Current weight-space learning approaches impose impractical complexity barriers: graph neural network methods require constructing computational graphs from weight tensors and 50+ hours of implementation, while forward-pass-based methods necessitate model instantiation and GPU resources. Both approaches prevent transparent, efficient analysis at scale.

This work demonstrates that architecture families impose structural constraints that manifest as checkpoint-observable fingerprints accessible through simple statistical features. CNNs use BatchNorm for spatial normalization of image data and allocate parameters predominantly to convolutional kernels. Transformers use LayerNorm for token-wise normalization and allocate parameters to large linear projection layers. These are not training conventions---they reflect fundamental data processing paradigms directly visible in checkpoint structure through layer names and tensor shapes.

We extract two features from PyTorch state\_dict: (1) normalization layer type counts (BatchNorm, LayerNorm, GroupNorm) via regex matching on layer names, and (2) parameter-mass ratio R = conv\_params / (conv\_params + linear\_params) via tensor shape inspection. A logistic regression classifier trained on these features achieves 88.89\% macro-averaged accuracy on 18 held-out TIMM models, exceeding the 80\% threshold by 8.89 percentage points. Edge case validation on non-standard architectures maintains 83.3\% accuracy with 1.7\% degradation.

Three mechanistic validations confirm feature discriminative power: (1) CNNs exhibit 0\% violation rate for BatchNorm usage, Transformers exhibit 14.29\% violation for LayerNorm (both within $\leq 15$\% threshold), confirming normalization choice reflects architectural paradigm. (2) Parameter-mass ratio exhibits Cohen's d = 3.202 (p < 0.001) between CNN and Transformer families, demonstrating fundamentally distinct allocation strategies. (3) Checkpoint-only extraction completes in 61 seconds with 0 MB GPU usage, enabling practical model zoo analysis.

The method exhibits perfect scale invariance: coefficient of variation = 0.00 across ResNet-{18,34,50,101,152}, spanning $5\times$ parameter range from 11M to 60M. This property enables family classification independent of model size without requiring size as auxiliary feature.

The principal limitation is complete failure on NormFree networks (0\% accuracy on NFNet) that replace normalization layers with scaled weight standardization. With absent normalization, classification relies solely on parameter-mass ratio, which is insufficient for these paradigms. The limitation affects rare architectures (VGG-16 historical, NFNets niche) and has clear mechanistic cause.

This work challenges the assumption that weight-space learning requires complex neural representations. When targeting specific classification tasks aligned with hand-crafted features guided by mechanistic understanding, simple statistical approaches achieve comparable accuracy with $100\times$ faster extraction and full interpretability.
